\documentclass{article}

\usepackage{graphicx}
\usepackage{amsmath}
\usepackage{amssymb}
\usepackage{algorithm}
\usepackage{algpseudocode}

\usepackage{listings}   % Para inserir blocos de código
\usepackage{changepage} % Para usar o ambiente adjustwidth

% (Opcional) Configurações básicas para deixar o código C mais bonito
\usepackage{xcolor} % Necessário se quiser adicionar cores ao código
\lstset{
	tabsize=4,
	basicstyle=\ttfamily\small,
	keywordstyle=\color{blue},
	stringstyle=\color{red},
	commentstyle=\color{green!60!black},
	showstringspaces=false,
	breaklines=true
}

\title{Coloração de Grafos}
\author{Bruno Henrique Duarte\\João Víctor Paiva Nunes Pereira}
\date{Abril 2026}

\begin{document}

\maketitle

\section{O Problema de Coloração de Grafos}

Grafos são estruturas matemáticas compostas por um conjunto de vértices $V$ e um conjunto de arestas $E$, geralmente denotados por $G=(V,E)$. Na Ciência da Computação (CC), problemas que podem ser abordados através de grafos abundam. Dentre esses problemas, a coloração de grafos merece destaque.

Uma $\lambda$-coloração de um grafo $G=(V,E)$ é uma função $f: V\rightarrow \{1,2,\dots,\lambda\}$ que mapeia cada vértice a uma cor (aqui representada por um número de 1 a $\lambda$). Uma coloração correta de um grafo é tal que para todo $x,y\in V$ tais que $y\in\Gamma(x)$ temos que $f(x)\not=f(y)$, onde $\Gamma: V\rightarrow 2^V$ retorna os vértices adjacentes - vizinhos - a um vértice dado.

Dada essa definição, dizemos que o número cromático de um grafo, denotado por $\chi(G)$, é a menor quantidade de cores necessária para colorir um grafo. Descobrir $\chi(G)$ é o problema de coloração de grafos.

\subsection{Problemas Equivalentes}

Diversos problemas práticos são equivalente ao problema de coloração de grafos, alguns deles:

\begin{enumerate}
	\item \textit{Resolução de um jogo de Sudoku} (Herzberg; Murty, 2006);
	
	\item \textit{Alocação de registradores em compiladores} (Aho \textit{et al.}, 2007, p. 353);
	
	\item \textit{Escalonamento}. Clássico problema geralmente associado à alocação de turmas, professores e alunos.
\end{enumerate}

Esses três problemas demonstram como o problema de coloração é um problema relevante na CC. Agora, vamos analisar os algoritmos que tentam resolver esse problema. Nesse sentido, é importante destacar que representamos grafos como lista de adjacências.

\section{Código}

\subsection{Código de Leitura do Grafo}

\subsubsection{arquivo.h}

\vspace{0.5em}
\begin{adjustwidth}{}{}
\begin{lstlisting}[language=C]
#ifndef _ARQUIVO_H
#define _ARQUIVO_H

#include "coloracao.h"

#define NOME_ARQUIVO "grafo_25nos.csv"
#define QTD_VERTICES 25
#define BUFFER_LEITURA 1 + 2*QTD_VERTICES

Grafo* lerArquivo();

#endif
\end{lstlisting}
\end{adjustwidth}
\vspace{0.5em}

\subsubsection{arquivo.c}

\vspace{0.5em}
\begin{adjustwidth}{}{}
\begin{lstlisting}[language=C]
#include <stdio.h>
#include <stdlib.h>
#include "arquivo.h"

Grafo* lerArquivo() {
	
	Grafo *g = criaGrafo(QTD_VERTICES);
	
	FILE *file = fopen(NOME_ARQUIVO, "r");
	if (!file) {
		#ifdef _WIN32
		SetConsoleOutputCP(CP_UTF8);
		#endif
		perror("\nNao foi possivel ler o arquivo.\n");
		exit(1);
	}
	
	char linha[BUFFER_LEITURA+1];
	linha[BUFFER_LEITURA] = '\0';
	
	int v = 0;
	while (v < QTD_VERTICES && fgets(linha, BUFFER_LEITURA + 1, file) != NULL) {
		
		int vAdj = 0;
		while (vAdj < QTD_VERTICES) {
			
			if (linha[2*vAdj] == '1') {
				Vertice *novaAdj = malloc(sizeof(Vertice));
				novaAdj->indice = vAdj;
				novaAdj->proximo = g->listaAdj[v].head;
				g->listaAdj[v].head = novaAdj;
			}
			vAdj++;
		}
		
		v++;
	}
	
	fclose(file);
	
	return g;
}
\end{lstlisting}
\end{adjustwidth}
\vspace{0.5em}

\subsection{Código de Coloração do Grafo}

\subsubsection{coloracao.h}

\vspace{0.5em}
\begin{adjustwidth}{}{}
\begin{lstlisting}[language=C]
#ifndef _COLORACAO_H
#define _COLORACAO_H

#ifdef _WIN32
#include <windows.h>
#endif

#define MIN_VERTICES 5
#define RANGE_ADD_VERTICES 15

typedef struct Vertice {
	int indice;
	struct Vertice* proximo;
} Vertice;

typedef struct Adjacencias {
	int cor;
	Vertice* head;
} Adjacencias;

typedef struct Grafo {
	int numVertices;
	int verticesColoridos;
	Adjacencias* listaAdj;
} Grafo;

void printaGrafo(Grafo* grafo, int printaCor);
void resetaCores(Grafo* grafo);

// Algoritmo de coloracao sobre o grafo
int coloreBacktracking(Grafo* grafo);
int coloreGuloso(Grafo *grafo);

// Verificacoes sobre o grafo
int temVerticeSemCor(Grafo *g);
int ehSeguro(Grafo* grafo, int v, int corTestada);
int estaBemColorido(Grafo *g);

// Cria ou destroi o grafo
Grafo* criaGrafo(int n);
void criaArestasRandom(Grafo* grafo);
Vertice* escolheVizinhos(Grafo* grafo, int i, int qtdAleat);
void freeGrafo(Grafo *g);

#endif
\end{lstlisting}
\end{adjustwidth}
\vspace{0.5em}

\subsubsection{coloracao.c}

\vspace{0.5em}
\begin{adjustwidth}{}{}
\begin{lstlisting}[language=C]
#include <stdio.h>
#include <stdlib.h>
#include <time.h>

#include "coloracao.h"
#include "arquivo.h"

void resetaCores(Grafo* grafo) {
	for(int i = 0; i < grafo->numVertices; i++) {
		grafo->listaAdj[i].cor = 0;
	}
	grafo->verticesColoridos = 0;
}

int temVerticeSemCor(Grafo *g) {
	return g->verticesColoridos != g->numVertices;
}

int ehSeguro(Grafo* grafo, int v, int corTestada) {
	Vertice* vizinho = grafo->listaAdj[v].head;
	
	while(vizinho != NULL) {
		if(grafo->listaAdj[vizinho->indice].cor == corTestada)
		return 0;
		
		vizinho = vizinho->proximo;
	}
	return 1;
}

int coloreBacktrackingAux(Grafo* grafo, int verticeAtual, int qtdCores) {
	if(!temVerticeSemCor(grafo)) {
		return 1;
	}
	
	for(int c = 1; c <= qtdCores; c++) {
		if(ehSeguro(grafo, verticeAtual, c)) {
			grafo->listaAdj[verticeAtual].cor = c;
			grafo->verticesColoridos++;
			
			if(coloreBacktrackingAux(grafo, verticeAtual+1, qtdCores)) {
				return 1;
			} else {
				grafo->listaAdj[verticeAtual].cor = 0;
				grafo->verticesColoridos--;
			}
		}
	}
	
	return 0;
}

int coloreBacktracking(Grafo* grafo) {
	for(int m = 1; m <= grafo->numVertices; m++) {
		resetaCores(grafo);
		
		if(coloreBacktrackingAux(grafo, 0, m)) {
			return m;
		}
	}
	return grafo->numVertices;
}

void coloreGulosoAuxiliar(Grafo *grafo, int verticeAtual, int *paletaCores) {
	
	Vertice *vAdj = grafo->listaAdj[verticeAtual].head;
	while (vAdj) {
		int corVizinho = grafo->listaAdj[vAdj->indice].cor;
		if (corVizinho != 0) paletaCores[corVizinho-1] = 0;
		vAdj = vAdj->proximo;
	}
	
	int corEscolhida = 1;
	for (; corEscolhida < grafo->numVertices; corEscolhida++)
	if (paletaCores[corEscolhida-1])
	break;
	
	grafo->listaAdj[verticeAtual].cor = corEscolhida;
	
	vAdj = grafo->listaAdj[verticeAtual].head;
	while (vAdj) {
		int corVizinho = grafo->listaAdj[vAdj->indice].cor;
		if (corVizinho != 0) paletaCores[corVizinho-1] = 1;
		vAdj = vAdj->proximo;
	}
	
	grafo->verticesColoridos++;
	
	Vertice *proxVertice = grafo->listaAdj[verticeAtual].head;
	while (temVerticeSemCor(grafo) && proxVertice != NULL) {
		if (grafo->listaAdj[proxVertice->indice].cor == 0)
		coloreGulosoAuxiliar(grafo, proxVertice->indice, paletaCores);
		proxVertice = proxVertice->proximo;
	}
}

int coloreGuloso(Grafo *grafo) {
	
	int *paletaCores = malloc(sizeof(int) * grafo->numVertices);
	for (int i = 0; i < grafo->numVertices; i++)
	paletaCores[i] = 1;
	
	for (int i = 0; i < grafo->numVertices && temVerticeSemCor(grafo); i++)
	if (grafo->listaAdj[i].cor == 0)
	coloreGulosoAuxiliar(grafo, i, paletaCores);
	
	int numCores = 0;
	
	for (int i = 0; i < grafo->numVertices && numCores < grafo->numVertices; i++) {
		if (paletaCores[grafo->listaAdj[i].cor-1]) {
			numCores++;
			paletaCores[grafo->listaAdj[i].cor-1] = 0;
		}
	}
	
	return numCores;
}

void printaGrafo(Grafo* grafo, int printaCor) {
	printf("Lista de Adjacencias: \n");
	
	for(int i = 0; i < grafo->numVertices; i++) {
		
		if(printaCor) {
			printf("[ (COR %d) %d:", grafo->listaAdj[i].cor, i);
		} else {
			printf("[ %d:", i);    
		}
		
		Vertice* aux = grafo->listaAdj[i].head;
		
		while(aux != NULL) {
			printf(" -> %d", aux->indice);
			aux = aux->proximo;
		}
		
		printf(" ]\n");
	}
}

// Fisher-Yates (auxiliar de criaArestasRandom)
// Nao ha auto-relacionamento nem vertices sem adjacencias
Vertice* escolheVizinhos(Grafo* grafo, int i, int qtdAleat) {
	if (qtdAleat == 0) return NULL; 
	
	int* baralho = malloc((grafo->numVertices - 1) * sizeof(int));
	int pos = 0;
	
	for(int k = 0; k < grafo->numVertices; k++) {
		if(k != i) {
			baralho[pos] = k;
			pos++;
		}
	}
	
	for (int k = pos - 1; k > 0; k--) {
		int r = rand() % (k + 1); 
		int temp = baralho[k];
		baralho[k] = baralho[r];
		baralho[r] = temp;
	}
	
	Vertice* headListaNova = NULL;
	
	for(int k = 0; k < qtdAleat; k++) { 
		int vizinhoEscolhido = baralho[k];
		
		Vertice* novoVertice = malloc(sizeof(Vertice));
		novoVertice->indice = vizinhoEscolhido;
		novoVertice->proximo = headListaNova;
		headListaNova = novoVertice;
	}
	
	free(baralho);
	return headListaNova; 
}

void criaArestasRandom(Grafo* grafo) {
	int qtdAleat;
	
	for(int i = 0; i < grafo->numVertices; i++) {
		qtdAleat = (rand() % (grafo->numVertices - 1)) + 1;
		grafo->listaAdj[i].head = escolheVizinhos(grafo, i, qtdAleat);
	}
	
	// Garantir a simetria do grafo
	for (int i = 0; i < grafo->numVertices; i++) {
		
		Vertice *vAdj_i = grafo->listaAdj[i].head;
		while (vAdj_i != NULL) {
			
			Vertice *vAdj_j = grafo->listaAdj[vAdj_i->indice].head;
			while (vAdj_j != NULL) {
				if (vAdj_j->indice == i)
				break;
				vAdj_j = vAdj_j->proximo;
			}
			
			if (vAdj_j == NULL) {
				Vertice *novoVertice = malloc(sizeof(Vertice));
				novoVertice->indice = i;
				novoVertice->proximo = grafo->listaAdj[vAdj_i->indice].head;
				grafo->listaAdj[vAdj_i->indice].head = novoVertice;
			}
			
			vAdj_i = vAdj_i->proximo;
		}
	}
}

Grafo* criaGrafo(int n) {
	Grafo* grafo = malloc(sizeof(Grafo));
	grafo->numVertices = n;
	grafo->verticesColoridos = 0;
	grafo->listaAdj = malloc(sizeof(Adjacencias)*n);
	
	for(int i = 0; i < n; i++) {
		grafo->listaAdj[i].head = NULL;
		grafo->listaAdj[i].cor = 0;
	}
	
	return grafo;
}

int estaBemColorido(Grafo *g) {
	for (int i = 0; i < g->numVertices; i++) {
		Vertice *vAdj = g->listaAdj[i].head;
		while (vAdj) {
			if (g->listaAdj[vAdj->indice].cor == g->listaAdj[i].cor) {
				printf("\n[%i -> %i]\n", i, vAdj->indice);
				return 0;
			}
			vAdj = vAdj->proximo;
		}
	}
	return 1;
}

void freeGrafo(Grafo *g) {
	
	if (!g || g->numVertices == 0 || !g->listaAdj) return;
	
	for (int i = 0; i < g->numVertices; i++) {
		Vertice *vAdj = g->listaAdj[i].head;
		while (vAdj) {
			Vertice *aux = vAdj;
			vAdj = vAdj->proximo;
			free(aux);
		}
	}
	free(g->listaAdj);
	g->numVertices = 0;
	g->verticesColoridos = 0;
}

int main(int argc, char *argv[]) {
	
	srand(time(NULL));
	
	#ifdef _WIN32
	SetConsoleOutputCP(CP_UTF8);
	#endif
	
	Grafo *grafo;
	
	if (argc == 2 && argv[1][0] == '1') {
		grafo = lerArquivo();
	} else {
		grafo = criaGrafo(MIN_VERTICES + (rand() % (RANGE_ADD_VERTICES + 1)));
		criaArestasRandom(grafo);
	}
	
	printaGrafo(grafo, 0);
	int numCores;
	
	//////////////////////////////////
	
	numCores = coloreBacktracking(grafo);
	
	printf("\n-----------------------------------\n");
	printf("Backtracking: Grafo colorido com %d cores\n", numCores);
	printf("-----------------------------------\n\n");
	
	printaGrafo(grafo, 1);
	printf("\nEsta corretamente colorido? %s\n", estaBemColorido(grafo) ? "sim" : "nao");
	
	//////////////////////////////////
	
	resetaCores(grafo);
	
	//////////////////////////////////
	
	numCores = coloreGuloso(grafo);
	
	printf("\n-----------------------------------\n");
	printf("Guloso: Grafo colorido com %d cores\n", numCores);
	printf("-----------------------------------\n\n");
	
	printaGrafo(grafo, 1);
	printf("\nEsta corretamente colorido? %s\n", estaBemColorido(grafo) ? "sim" : "nao");
	
	//////////////////////////////////
	
	freeGrafo(grafo);
	
	return 0;
}
\end{lstlisting}
\end{adjustwidth}
\vspace{0.5em}

\section{Coloração via Backtracking}

Vamos iniciar pelo algoritmo força bruta para resolver o problema.

\vspace{0.5em}
\begin{adjustwidth}{2.5em}{}
\begin{lstlisting}[language=C]
int coloreBacktrackingAux(Grafo* grafo, int verticeAtual, int qtdCores) {
	if(!temVerticeSemCor(grafo)) {
		return 1;
	}
	
	for(int c = 1; c <= qtdCores; c++) {
		if(ehSeguro(grafo, verticeAtual, c)) {
			grafo->listaAdj[verticeAtual].cor = c;
			grafo->verticesColoridos++;
			
			if(coloreBacktrackingAux(grafo, verticeAtual+1, qtdCores)) {
				return 1;
			} else {
				grafo->listaAdj[verticeAtual].cor = 0;
				grafo->verticesColoridos--;
			}
		}
	}
	
	return 0;
}

int coloreBacktracking(Grafo* grafo) {
	for(int m = 1; m <= grafo->numVertices; m++) {
		resetaCores(grafo);
		
		if(coloreBacktrackingAux(grafo, 0, m)) {
			return m;
		}
	}
	return grafo->numVertices;
}
\end{lstlisting}
\end{adjustwidth}
\vspace{0.5em}

A função que é chamada inicialmente é a \texttt{coloreBacktracking}. Nesta função, há um laço que tenta colorir o grafo com uma quantidade de cores $m$, onde $1\le m\le|V|$. Para cada valor de $m$, o laço chama a função auxiliar \texttt{coloreBacktrackingAux}, que é a função que de fato tenta colorir o grafo com a quantidade de cores dada e com o vértice inicial passado como argumento (sempre o vértice 0).

Na função \texttt{coloreBacktrackingAux}, é realizado um laço sobre as cores. Se o vértice atual puder ser colorido com a cor (verificação feita pela função \texttt{ehSeguro}), colore-se ele e chama-se recursivamente a função para o vértice de índice seguinte. Quando a recursão acabar, caso seja retornado 1, significa que é possível colorir o grafo com a quantidade de cores dada; caso seja retornado 0, significa que não foi possível colorir algum vértice com nenhuma das cores disponíveis.

A complexidade de espaço da função \texttt{coloreGuloso} é $O(V)$, enquanto a complexidade de tempo é $O(V^V)$. Para enxergar a complexidade de espaço, basta verificar que a recursão de \texttt{coloreBacktrackingAux} sempre para ao ser chamada pela $V+1$ vez e retorna até a chamada inicial realizada na função \texttt{coloreGuloso}.

Quanto à complexidade de tempo, temos que olhar para o laço principal de \texttt{coloreBacktrackingAux}. Seja $m$ a quantidade de cores, o laço sempre executa $m$ vezes e chama a função \texttt{ehSeguro}, a qual tem complexidade $O(E)$, mas que para um grafo denso pode ser $O(V)$. Caso seja seguro colorir o vértice, a função realiza a recursão. Neste ponto é onde a árvore de execução se ramifica e conclui-se que a complexidade de tempo da função \texttt{coloreBacktrackingAux} é $O(V\cdot m^V)$. Voltando à função \texttt{coloreBacktracking}, o laço executa \texttt{resetaCores}, que tem complexidade de tempo $O(V)$, e a função auxiliar. Logo, a complexidade é dada por:

\[\sum_{m=1}^{V} [O(V) + O(V\cdot m^V)] = O(V^V).\]

Evidentemente, essa é uma complexidade muito ruim. Isso se deve ao fato de que a coloração de grafos é um problema NP-difícil, ou seja, não se conhece algoritmo em tempo polinomial que resolva otimamente esse problema.

\section{Coloração via Técnica Gulosa}

Diante do fato do problema ser NP-difícil, procuram-se outras estratégias para realizar a coloração que sejam muito mais eficientes. Em troca, essas abordagens não encontram a soluçam ótima. O algoritmo a seguir utiliza uma técnica gulosa para colorir o grafo.

\vspace{0.5em}
\begin{adjustwidth}{}{}
\begin{lstlisting}[language=C]
void coloreGulosoAuxiliar(Grafo *grafo, int verticeAtual, int *paletaCores) {
	
	Vertice *vAdj = grafo->listaAdj[verticeAtual].head;
	while (vAdj) {
		int corVizinho = grafo->listaAdj[vAdj->indice].cor;
		if (corVizinho != 0) paletaCores[corVizinho-1] = 0;
		vAdj = vAdj->proximo;
	}
	
	int corEscolhida = 1;
	for (; corEscolhida < grafo->numVertices; corEscolhida++)
	if (paletaCores[corEscolhida-1]) break;
	
	grafo->listaAdj[verticeAtual].cor = corEscolhida;
	
	vAdj = grafo->listaAdj[verticeAtual].head;
	while (vAdj) {
		int corVizinho = grafo->listaAdj[vAdj->indice].cor;
		if (corVizinho != 0) paletaCores[corVizinho-1] = 1;
		vAdj = vAdj->proximo;
	}
	
	grafo->verticesColoridos++;
	
	Vertice *proxVertice = grafo->listaAdj[verticeAtual].head;
	while (temVerticeSemCor(grafo) && proxVertice != NULL) {
		if (grafo->listaAdj[proxVertice->indice].cor == 0)
		coloreGulosoAuxiliar(grafo, proxVertice->indice, paletaCores);
		proxVertice = proxVertice->proximo;
	}
}

int coloreGuloso(Grafo *grafo) {
	
	int *paletaCores = malloc(sizeof(int) * grafo->numVertices);
	for (int i = 0; i < grafo->numVertices; i++)
	paletaCores[i] = 1;
	
	for (int i = 0; i < grafo->numVertices && temVerticeSemCor(grafo); i++)
		if (grafo->listaAdj[i].cor == 0)
			coloreGulosoAuxiliar(grafo, i, paletaCores);
	
	int numCores = 0;
	
	for (int i = 0; i < grafo->numVertices && numCores < grafo->numVertices; i++) {
		if (paletaCores[grafo->listaAdj[i].cor-1]) {
			numCores++;
			paletaCores[grafo->listaAdj[i].cor-1] = 0;
		}
	}
	
	return numCores;
}
\end{lstlisting}
\end{adjustwidth}
\vspace{0.5em}

Inicialmente, chama-se o procedimento \texttt{coloreGuloso}, o qual cria um vetor de inteiros de tamanho $V$. Esse vetor, na verdade, opera como um vetor booleano, em que cada posição indica se uma cor pode ser utilizada para pintar o vértice ou não. Nesse procedimento, inicializa-se o vetor com 1s. Na sequência, o procedimento chama a função auxiliar \texttt{coloreGulosoAuxiliar}, passando como argumento o vértice inicial para começar a coloração e o vetor booleano.

Na função \texttt{coloreGulosoAuxiliar}, primeiro varre-se as adjacências do vértice atual e marca-se com 0 a posição das cores utilizadas nos seus vizinhos. A seguir, encontra-se a primeira cor não marcada para colorir o vértice atual. Então, reseta-se o vetor booleano e chama-se a função recursivamente para algum vértice vizinho ainda não colorido. Caso não haja tal vértice, a função retorna para a função principal.

Evidentemente, a complexidade de espaço dessa abordagem gulosa é $O(V)$, para o vetor booleano. Já a complexidade de tempo é $O(V+E)$, pois a complexidade de \texttt{coloreGulosoAuxiliar} é $O(V+E)$ e essa função só é chamada enquanto houver vértices que não foram coloridos.

A técnica gulosa é empregada no procedimento \texttt{coloreGulosoAuxiliar}, onde procura-se tentar colorir primeiro os vértices vizinhos. Eventualmente, todos os vértices terão sido coloridos. O único problema é que isso não garante que encontraremos o número cromático do grafo. Todavia, para a maioria dos problemas uma aproximação boa é o suficiente.

\section{Resultados}

Como esperado, o algoritmo que utiliza backtracking levou mais de 5 minutos para encontrar o número cromático do grafo dado pelo professor. O número cromático encontrado foi 8. Para grafos pequenos, o backtracking pode até ser uma abordagem viável, porém continuará sendo muito lento.

Por outro lado, o algoritmo guloso executou extremamente rápido. Diante disso, foi possível realizar diversas execuções com vértices iniciais diferentes para analisar se a quantidade de cores utilizadas para colorir o grafo mudaria.

Dessas execuções, obtemos:

\begin{itemize}
	\item Para 5 vértices iniciais o grafo foi colorido com 10 cores;
	\item Para 18 vértices iniciais o grafo foi colorido com 9 cores;
	\item Para 2 vértices iniciais o grafo foi colorido com 8 cores, a solução ótima.
\end{itemize}

Desse modo, podemos concluir que a maioria das execuções retornou um número muito próximo à solução ótima, que o vértice inicial de coloração interfere diretamente na quantidade de cores que serão utilizadas e que o algoritmo guloso é uma excelente aproximação - e às vezes uma solução - para o problema de coloração de grafos.

\section{Referências}

AHO, A. V.; LAM, M. S.; SETHI, R.; ULLMAN, J. D. Compilers: principles, techniques, and tools. 2. ed. Boston: Pearson Addison-Wesley, 2006.

\noindent HERZBERG, A. M.; MURTY, M. R. Sudoku squares and chromatic polynomials. Notices of the American Mathematical Society, Providence, v. 54, n. 6, p. 708-717, jun. 2007.

\end{document}